\section{Target System}\label{sec:target}
%% Oscar Gustafsson

In the target system, a ring of $N$ friendly drones patrols at radius $R_p$
around a protected military camp. Each drone carries a camera on a gimbal
that scans the sky outside the ring; which camera, and whether a second one
is added, is chosen in Section~\ref{sec:camera-choice}. A hostile drone
approaches from outside at speed $v_h$. The nearest friendly drone sees it as
a small dark object against the sky. A fast model on the drone flags the
object, and the drone starts flying towards it at once, without waiting for
the base station. Throughout its flight, it sends an image crop of the object
and the direction to it every $\Delta t$ seconds. The base station acts on
these (Section~\ref{sec:base-station}): a stronger model classifies the
object, and if it is a drone, the base station confirms the intercept and
sends a second drone so that the target can be triangulated; if it is not,
the base station calls the drone back. Since missing a hostile drone is worse
than a detour, the drone is only called back when the base station is
confident that the object is not a drone. The interceptor acquires and locks
on to the target, keeps it centred in its gimbal camera, and guides itself to
the intercept. All
numbers below follow from the assumptions in Section~\ref{sec:assumptions} and
Section~\ref{sec:target-assumptions}; when an assumption changes, the derived
values change with it.

\subsection{Additional Assumptions}
\label{sec:target-assumptions}

The target system adds components and constraints that the baseline does not
have. Table~\ref{tab:target-assumptions} lists their assumed values, on top of
those in Table~\ref{tab:assumptions}.

\begin{table}[H]
\centering
\small
\begin{tabular}{|l|p{0.25\textwidth}|r|p{0.27\textwidth}|}
\hline
\textbf{Symbol} & \textbf{Quantity} & \textbf{Value} & \textbf{Note} \\
\hline
$\Theta_s$ & Sector scanned by the gimbal camera & 180$^\circ$ & Constraint: outward half only, since threats come from outside \\
$\Delta t$ & Interval between image crops sent to the base station & 0.2\,s & Addition: 5 crops per second \\
$t_c$ & Image sequence needed to classify & 1.0\,s & Addition: base-station classifier, $t_c/\Delta t = 5$ crops \\
$t_l$ & Radio link and inference latency & 0.1\,s & Addition: drone to base station and back \\
$p_{\min}$ & Smallest detectable size & 4\,px & Analysis: dark object against the sky \\
$t_{\mathrm{lock}}$ & Lock-on time that counts as a successful interception & 3\,s & Evaluation: alternative to a hit (Section~\ref{sec:intercept}) \\
$d_{\mathrm{lock}}$ & Largest distance for a lock-on & 20\,m & Evaluation: close enough to engage the target \\
\hline
\end{tabular}
\caption{Additional assumed values of the target system.}
\label{tab:target-assumptions}
\end{table}

By Johnson's criteria, detecting that an object is present needs about one
line pair, that is two pixels, across it \cite{johnson1958analysis};
$p_{\min} = 4$\,px doubles that as a margin. For the analysis below, a camera
is also assumed to detect the target as soon as it is at least $p_{\min}$
pixels wide, images have a 16:9 aspect ratio, and
the ring is treated as locally straight where the hostile drone crosses it.
The camera's field of view $\theta$ and image width $W$ are left open; they
are the design variables of Section~\ref{sec:camera-choice}.

\subsection{Derived Quantities}
\label{sec:derived}

\textbf{Drone spacing (sparsity).} With $N$ drones evenly spaced on the ring,
the distance between neighbours is
\begin{equation}
  d = \frac{2\pi R_p}{N},
  \label{eq:spacing}
\end{equation}
and a hostile drone that crosses the ring midway between two neighbours passes
at a lateral distance of at most
\begin{equation}
  y_{\max} = \frac{d}{2} = \frac{\pi R_p}{N}
  \label{eq:ymax}
\end{equation}
from the nearest friendly drone.

\textbf{Detection range.} By the pinhole camera model
\cite{hartley2004multiple}, a camera of width $W$ pixels and horizontal field
of view $\theta$ has a focal length, in pixels, of
\begin{equation}
  f = \frac{W/2}{\tan(\theta/2)}.
  \label{eq:focal}
\end{equation}
A target of size $s$ at range $r$ covers $p = f s / r$ pixels, so the
detection range is
\begin{equation}
  r_{\mathrm{det}} = \frac{f\,s}{p_{\min}}.
  \label{eq:rdet}
\end{equation}

\textbf{Coverage without gaps.} Every hostile drone that crosses the ring is
seen if $r_{\mathrm{det}} \geq y_{\max}$, which by Equation~\eqref{eq:ymax}
requires
\begin{equation}
  N \geq \left\lceil \frac{\pi R_p}{r_{\mathrm{det}}} \right\rceil.
  \label{eq:ncoverage}
\end{equation}

\textbf{Scanning.} A camera with image size $W \times H$ has the vertical field
of view
\begin{equation}
  \theta_v = 2\arctan\!\left(\tan\frac{\theta}{2}\cdot\frac{H}{W}\right).
  \label{eq:vfov}
\end{equation}
At range $r_{\mathrm{det}}$, hostile altitudes between $h_{h,\min}$ and
$h_{h,\max}$ span the elevation band
$\Delta\varepsilon = \arctan\frac{h_{h,\max}-h}{r_{\mathrm{det}}}
 - \arctan\frac{h_{h,\min}-h}{r_{\mathrm{det}}}$.
Covering the outward sector $\Theta_s$ then takes
$n = \lceil \Theta_s/\theta \rceil \cdot \lceil \Delta\varepsilon/\theta_v \rceil$
camera positions, and one sweep takes
\begin{equation}
  T_{\mathrm{sweep}} = n\,t_g ,
  \label{eq:sweep}
\end{equation}
during which the hostile drone flies $v_h T_{\mathrm{sweep}}$.

\textbf{Time to depart.} In the worst case, the hostile drone enters the
detection range just after the camera has looked past it. The drone that
detects it starts flying as soon as its own model flags the object, so the
time from entering detection range to departure is
\begin{equation}
  t_d = T_{\mathrm{sweep}} + t_g .
  \label{eq:delay}
\end{equation}
The base station's decision arrives $t_c + t_l$ later, while the drone is
already on its way, so it does not delay the intercept; it only decides
whether the flight continues.

\textbf{Intercept feasibility.} Since $v_h > v_i$, the interceptor cannot catch
the hostile drone from behind; it must reach the hostile drone's path before
the hostile drone gets there. Conservatively, the meeting point is where the
hostile drone crosses the ring. A hostile drone first detected at range
$r_{\mathrm{det}}$ by a drone at lateral distance $y$ reaches the ring after
\begin{equation}
  t_h = \frac{\sqrt{r_{\mathrm{det}}^2 - y^2}}{v_h},
  \label{eq:th}
\end{equation}
and the interceptor reaches the hostile drone's path after
\begin{equation}
  t_i = t_d + \frac{y}{v_i} + \frac{v_i}{2a},
  \label{eq:ti}
\end{equation}
where the last term is the time lost while accelerating to $v_i$. The
intercept is feasible when $t_h \geq t_i$ at the worst case $y = y_{\max}$.

\textbf{Keeping the target in view.} A target that moves across the line of
sight at speed $v_\perp$ and range $r$ turns, as seen by the camera, at the
rate
\begin{equation}
  \omega = \frac{v_\perp}{r},
  \label{eq:angrate}
\end{equation}
and leaves the image after $(\theta/2)/\omega$ if the gimbal stops following
it.

\subsection{Camera Choice}
\label{sec:camera-choice}

The camera is a trade-off. A narrower field of view or a wider image raises
the focal length (Equation~\eqref{eq:focal}) and so the detection range, but a
narrower field of view needs more camera positions to scan the sector, which
lengthens the sweep and the time to decide. Table~\ref{tab:coverage} shows the
effect on the number of drones for a range of cameras, including the one the
simulated platform already has.

\begin{table}[H]
\centering
\small
\begin{tabular}{|r|r|r|r|r|r|r|}
\hline
$\theta$ & $W$ (px) & $r_{\mathrm{det}}$ & Positions $n$ & $T_{\mathrm{sweep}}$ & $N$, no gaps & $N$, feasible \\
\hline
115$^\circ$ & 1280 & 51\,m   & 4  & 2.0\,s & 62 & -- \\
115$^\circ$ & 3840 & 153\,m  & 2  & 1.0\,s & 21 & 65 \\
60$^\circ$  & 1920 & 208\,m  & 3  & 1.5\,s & 16 & 42 \\
30$^\circ$  & 1920 & 448\,m  & 6  & 3.0\,s & 8  & 16 \\
20$^\circ$  & 1920 & 681\,m  & 9  & 4.5\,s & 5  & 10 \\
10$^\circ$  & 1920 & 1372\,m & 18 & 9.0\,s & 3  & 5 \\
20$^\circ$  & 3840 & 1361\,m & 9  & 4.5\,s & 3  & 5 \\
\hline
\end{tabular}
\caption{Detection range, scan time and the number of drones needed, from
Equations~\eqref{eq:spacing}--\eqref{eq:ti}. The first row is the camera of
the simulated platform (Table~\ref{tab:assumptions}); -- means that no number
of drones makes the intercept feasible.}
\label{tab:coverage}
\end{table}

Three conclusions follow.

\textbf{The platform's camera is not enough.} With the camera of the simulated
platform, a 0.5\,m drone is detectable at only about 51\,m, and no number of
drones on a 1\,km ring makes the intercept feasible. The target system
therefore needs a camera with a longer focal length than the simulation's
default.

\textbf{Within the model, narrower is better, with diminishing returns.} At
1920\,px, going from 60$^\circ$ to 20$^\circ$ cuts the drones needed from 42
to 10; going further to 10$^\circ$ saves only five more. The model
does not limit how narrow the camera can be, so the limit must come from what
it leaves out:
\begin{itemize}
  \item \emph{Keeping the target in view.} By Equation~\eqref{eq:angrate}, a
        target crossing at 17\,m/s turns at about 19.5$^\circ$/s at 50\,m and
        49$^\circ$/s at 20\,m. A 20$^\circ$ camera loses it after about
        0.2\,s at 20\,m if the gimbal lags; a 115$^\circ$ camera keeps it for
        about 1.2\,s. A narrow camera is good for finding the target far away
        but poor for keeping it close in.
  \item \emph{Computing cost.} The detector's work grows with the number of
        pixels: a 1920$\times$1080 image has 2.25 times, and a
        3840$\times$2160 image 9 times, the pixels of the platform's
        1280$\times$720 image. This matters on board the drone and when many
        drones render cameras in the same simulation.
  \item \emph{Pointing and stabilisation.} The narrower the camera, the more
        precisely the gimbal must point and hold still; a pointing error of
        half the field of view already loses the target.
\end{itemize}

\textbf{Options to test.} These effects are not known well enough to fix the
camera now. The target system is therefore evaluated in simulation with the
following options, using the platform's camera as the reference:
\begin{enumerate}
  \item \emph{One narrower gimbal camera}, about 20--30$^\circ$ at 1920\,px:
        the simplest option that makes the intercept feasible
        (Table~\ref{tab:coverage}).
  \item \emph{A narrow and a wide camera on the same gimbal}, or a zoom
        camera, which is wide when close and narrow when far: the narrow view
        finds the target far away, and the wide view keeps it in view at
        close range, where Equation~\eqref{eq:angrate} shows the narrow
        camera struggles. The wide view also shows where to turn the gimbal
        so that the narrow camera sees an object (\emph{slew-to-cue}), as in
        master--slave surveillance systems where a wide camera directs a
        pan--tilt--zoom camera \cite{delbimbo2010exploiting}. In Gazebo, both
        variants can be built as two camera sensors on the gimbal.
  \item \emph{Several narrow cameras side by side} (a camera array): nine
        20$^\circ$ cameras cover the 180$^\circ$ sector at once, so no
        scanning is needed and $T_{\mathrm{sweep}}$ disappears. Once a target
        is detected, lock-on needs only the one camera that sees it. The gain
        is smaller than it seems, because most of the intercept time is the
        flight, not the search: with $t_d \approx 0$, the minimum number of
        drones falls from 10 to 9, and the margin with 12 drones grows from
        7\,s to 12\,s. Three costs come with it:
        \begin{itemize}
          \item \emph{Detection runs on every camera.} Until something is
                found, all cameras must be searched: nine 1920$\times$1080
                images are 20 times the pixels of the platform's camera per
                frame. Processing them in turn at 30 images per second keeps
                the computing cost of one camera and still gives a cycle of
                0.3\,s, which leaves the minimum at 9 drones.
          \item \emph{The array tilts with the drone.} Fixed to the drone, the
                cameras pitch with it (Equation~\eqref{eq:tilt}); with a
                vertical field of view of only 11.3$^\circ$, covering a pitch
                of $\pm15^\circ$ takes four rows, that is 36 cameras. The
                array must therefore be stabilised, for example on a gimbal,
                which then carries nine cameras instead of one.
          \item \emph{Weight, power and simulation.} Every camera adds weight,
                power and bandwidth on a real drone, and rendering time in
                Gazebo, multiplied by the number of drones.
        \end{itemize}
\end{enumerate}
The option that intercepts most often with the fewest drones is chosen. Until
then, the analysis uses 20$^\circ$ at 1920\,px as its example.

\textbf{Rejected: a wide camera fixed to the drone.} A multirotor tilts to
accelerate: holding its altitude while accelerating horizontally at $a$, it
pitches by
\begin{equation}
  \phi = \arctan\frac{a}{g},
  \label{eq:tilt}
\end{equation}
where $g$ is the gravitational acceleration; $a = 2$\,m/s$^2$ already gives
$\phi \approx 11.5^\circ$. It tilts further when it brakes, turns or flies into wind. A camera fixed to the
drone tilts with it, so a target at elevation $\varepsilon$ above the horizon
stays in view only if
\begin{equation}
  \varepsilon + \phi \leq \frac{\theta_v}{2}.
  \label{eq:tiltview}
\end{equation}
Even a 115$^\circ$ camera, with $\theta_v \approx 83^\circ$, keeps only
targets less than 11.5$^\circ$ above the horizon in view at a pitch of
30$^\circ$, and the image shifts every time the drone changes its attitude.
This happens exactly during the intercept, when the drone accelerates hardest
and the target is close and high in the image. A gimbal removes the drone's
pitch, so all cameras of the target system are mounted on the gimbal.

\textbf{Example: 20$^\circ$ at 1920\,px.} The camera's vertical field of view
is 11.3$^\circ$ (Equation~\eqref{eq:vfov}). At $r_{\mathrm{det}} = 681$\,m,
hostile altitudes of 20--120\,m appear between $-2.5^\circ$ and $+5.9^\circ$
in elevation, which fits within one row of images, so the outward sector takes
$n = 9$ positions, $T_{\mathrm{sweep}} = 4.5$\,s and $t_d = 5.0$\,s.
Table~\ref{tab:feasibility} shows the resulting margin.

\begin{table}[H]
\centering
\small
\begin{tabular}{|r|r|r|r|r|r|}
\hline
$N$ & $d$ & $y_{\max}$ & $t_h$ & $t_i$ & Margin $t_h - t_i$ \\
\hline
8  & 785\,m & 393\,m & 32.7\,s & 40.7\,s & $-8.0$\,s \\
9  & 698\,m & 349\,m & 34.4\,s & 37.1\,s & $-2.7$\,s \\
10 & 628\,m & 314\,m & 35.5\,s & 34.2\,s & 1.3\,s \\
12 & 524\,m & 262\,m & 37.0\,s & 29.8\,s & 7.1\,s \\
16 & 393\,m & 196\,m & 38.3\,s & 24.4\,s & 14.0\,s \\
20 & 314\,m & 157\,m & 39.0\,s & 21.1\,s & 17.9\,s \\
\hline
\end{tabular}
\caption{Worst-case intercept margin for different numbers of drones, with a
20$^\circ$ camera at 1920\,px, from Equations~\eqref{eq:spacing}--\eqref{eq:ti}.}
\label{tab:feasibility}
\end{table}

With this camera, coverage is gap-free from $N = 5$, but an intercept is only
feasible in the worst case from $N = 10$; about 12 drones leave a margin of
7\,s.

\subsection{Detection and Interception}
\label{sec:target-approach}

\textbf{Detection.} Detection uses the drones' own cameras. The gimbal
cameras scan the sky outside the ring; how far they see, and how many drones
that requires, depends on the camera chosen in
Section~\ref{sec:camera-choice}. If a wide camera is added on the same
gimbal, it shows where to turn the gimbal for the narrow camera to see an
object (\emph{slew-to-cue}) and keeps the target in
view at close range. Small, distant drones are the hardest to detect, and
easily confused with birds, especially from a moving camera
\cite{coluccia2021drone}. Classification uses two tiers: a fast model on the
drone, tuned so that it rarely misses a drone, and a stronger model at the
base station that uses short image sequences and several drones' views.
Missing a hostile drone is treated as the worst outcome, so both tiers favour
recall; a false alarm only costs a detour, during which the target's image
grows and the classification improves.

\textbf{Interception.} A planning-based controller serves as the baseline: the
interceptor flies to the predicted meeting point on the hostile drone's path,
then steers with proportional navigation once the camera has locked on
\cite{zarchan2019tactical}. A controller based on reinforcement learning is
trained in a fast, simplified simulator; such controllers have been trained
to intercept using only the direction to the target
\cite{gaudet2019reinforcement}. It runs through the same drone interface, so
that the two can be compared on identical scenarios.

\subsection{System Overview}

Figure~\ref{fig:final-system} shows the target system. The base station
classifies detections and decides which drone does what; each drone runs its
own detection, lock-on and controller, so it can finish an intercept even if
the radio link is lost.

\begin{figure}[H]
\centering
\small
\begin{verbatim}
 BASE STATION
   tracks --> classifier (strong) --> assignment
     ^                                   |
     | image crops + bearings            | roles: patrol /
     |                                   v investigate / intercept
 EACH DRONE
   gimbal camera(s), scanning (narrow, + optional wide) --> fast detector
     --> lock-on (detect -> acquire -> locked)
     --> controller (planning or RL) --> PX4 autopilot --> Gazebo
\end{verbatim}
\caption{The target system.}
\label{fig:final-system}
\end{figure}
